\documentclass[11pt]{article}
\usepackage[a4paper,margin=26mm]{geometry}
\usepackage[T1]{fontenc}
\usepackage{amsmath,amssymb,amsthm,booktabs,array}
\usepackage[hidelinks]{hyperref}
\setlength{\emergencystretch}{1em}
\newtheorem{theorem}{Theorem}
\newtheorem{proposition}[theorem]{Proposition}
\newtheorem{corollary}[theorem]{Corollary}
\newcommand{\Om}{\Omega}
\title{Boundary corrections and a length-three extension of\texorpdfstring{\\}{}
B\'ona--Maga--Richey's extremal-frequency analysis}
\author{Haobo Ma\\\small ChronoAI Pte Ltd, Singapore
\and Wenlin Zhang\\\small National University of Singapore\\\small ChronoAI Pte Ltd, Singapore}
\date{November 16, 2026}
\begin{document}
\maketitle
\begin{abstract}
For B\'ona, Maga and Richey's extremal-frequency question
(arXiv:2606.06655v2, revised September~9, 2026, Section~5.3, p.~19),
we extend the classical frequency calculations of Lekkerkerker (1951),
Chang (1986) and Kova\v{c}evi\'c (2018/2019) to classify the maximizers
at lengths one, two and three by exact counting and generating functions.
They are $0$, $00$, and the pair $100,001$, respectively: the first two
lengths contradict the unrestricted proposal, whereas length three agrees
with it. We prove the authors' reducible-word frequency identity with its
necessary boundary restriction, obtain all four complementary families,
and show that their candidate frequencies decrease strictly to $1/2$.
The comparison with all competitors at lengths at least four is separate.
\end{abstract}

\section{The question and its short-length answer}
For a nonempty binary word $w$ and an integer $n\geq1$, let
\[
 \Om_n^w=\{u\in\{0,1\}^n: w\text{ is not a contiguous factor of }u\},
 \qquad
 \rho_n^w=\frac{\sum_{u\in\Om_n^w}|u|_1}{n|\Om_n^w|}.
\]
Here $|u|_1$ counts the ones in $u$, and $\rho^w=\lim_{n\to\infty}\rho_n^w$
when the limit exists. These are Definition~2.4 of~\cite{BMR}.
The denominator is positive: for nonempty $w$, at least one constant binary
word avoids $w$ at every length.
The paper defines words for all lengths $k\geq1$ (Section~2, p.~3).
Its Section~5.3, p.~19, proposes that forbidding $10^{k-1}$ or its reversal
maximizes $\rho^w$ over all words of length $k$: ``forbidding the word
$10^{k-1}$ or its reversal should maximize the frequency $\rho$''.
Theorem~\ref{thm:boundary} supplies counterexamples at the first two lengths;
Theorem~\ref{thm:three} proves the proposed comparison at length three.

The following classification shows why a restriction on $k$ is necessary.
The golden-mean value already follows from Lekkerkerker's 1951 theorem
on Fibonacci representations~\cite[Stelling~B, p.~2]{L}: deleting the
forced initial $10$ identifies representations of length $n+2$ with
all $11$-avoiding words of length $n$, and changes the number of ones by $1$.
St.~John~\cite[Theorem~(2)]{SJ} gives a direct proof for unrestricted first
letters. B\'ona--Maga--Richey~\cite[Example~2.10, Eq.~(2.14)]{BMR}
rederive the same value. We include its calculation for the extremal comparison.

\begin{theorem}\label{thm:boundary}
For one-letter forbidden words,
$\rho_n^0=1$ and $\rho_n^1=0$ for every $n\geq1$.
For two-letter forbidden words the limits are
\[
\begin{array}{c|cccc}
 w&00&01&10&11\\ \hline
 \rho^w&\displaystyle\frac12+\frac{1}{2\sqrt5}
 &\displaystyle\frac12&\displaystyle\frac12
 &\displaystyle\frac12-\frac{1}{2\sqrt5}.
\end{array}
\]
Consequently the unique maximizers at lengths one and two are $0$ and $00$,
respectively. The proposal on the full domain $k\geq1$ is false, and it
remains false if that domain is changed to $k\geq2$.
\end{theorem}

\begin{proof}
Avoiding $0$ leaves just $1^n$, and avoiding $1$ leaves just $0^n$.
A word avoiding $10$ has the unique form $0^a1^{n-a}$, $0\leq a\leq n$.
There are $n+1$ such words, containing $n(n+1)/2$ ones altogether;
thus $\rho_n^{10}=1/2$. Reversing words preserves both avoidance and the
number of ones, so $\rho_n^{01}=1/2$ as well.

For $00$, let $A_n(y)=\sum_{u\in\Om_n^{00}}y^{|u|_1}$, including
$A_0(y)=1$. An avoiding word is either empty, the single word $0$, or
starts with $1$ or $01$, followed by another avoiding word. Hence
\[
 \sum_{n\geq0}A_n(y)t^n=\frac{1+t}{1-yt-yt^2},\qquad
 A_1(y)=1+y,\quad A_n(y)=yA_{n-1}(y)+yA_{n-2}(y)\ (n\geq2).
\]
The characteristic roots are
$\lambda_\pm(y)=(y\pm\sqrt{y^2+4y})/2$.
For real $y$ near $1$, $\lambda_+(y)>|\lambda_-(y)|$ and
\[
 A_n(y)=C_+(y)\lambda_+(y)^n+C_-(y)\lambda_-(y)^n,
 \qquad C_+(1)=\frac{\varphi^2}{\sqrt5}>0,
\]
where $\varphi=(1+\sqrt5)/2$ and $C_\pm$ are differentiable near $1$.
This follows by solving the recurrence with the two displayed initial values.
The strict root gap persists on a neighborhood of $1$; differentiating
the explicit formula and dividing by $nA_n(1)$ therefore gives
\[
 \rho^{00}=\lim_{n\to\infty}\frac{A'_n(1)}{nA_n(1)}
 =\frac{\lambda'_+(1)}{\lambda_+(1)}.
\]
Differentiating $\lambda_+^2-y\lambda_+-y=0$ yields
$\lambda'_+(1)=(\varphi+1)/(2\varphi-1)$, and consequently
\[
 \rho^{00}=\frac{\varphi}{\sqrt5}
 =\frac12+\frac{1}{2\sqrt5}.
\]
Bit-complement is a bijection from $\Om_n^{00}$ to $\Om_n^{11}$ that
replaces $|u|_1$ by $n-|u|_1$, giving the remaining value.
The strict inequalities in the table prove uniqueness.
\end{proof}

\section{Why the reduction changes at length two}
Section~3.1, p.~8, of~\cite{BMR} observes that forbidding $10^m$ has the same
limiting letter frequency as forbidding $0^m$.
The following argument establishes that observation for $m\geq2$ and
identifies its exceptional boundary. Its frequency formula also gives the
candidate value in the longer-length extremal question.
The constant-word frequencies were computed by Chang~\cite[p.~179]{C}
and later by Baril--Kirgizov--Vajnovszki~\cite[Theorem~1]{BKV}.
The three-word identity below also follows from
Kova\v{c}evi\'c~\cite[Lemmas~1(b), 2 and Remark~4]{K}:
set the run parameters to $(d,k)=(0,m-1)$; his typical relative weight
is $f_m$, and his boundary observation allows either end run to be free.
Concentration of a variable in $[0,1]$ implies convergence of its mean.
The proof below gives a direct generating-function derivation.

\begin{proposition}\label{prop:reduction}
Let $m\geq2$, and let $r_m\in(0,1)$ be the unique positive solution of
$r_m+r_m^2+\cdots+r_m^m=1$. Then
\[
 \rho^{10^m}=\rho^{0^m1}=\rho^{0^m}
 =\frac{1}{\sum_{i=1}^m i r_m^i}.
\]
For $m=1$, instead, $\rho^{10}=\rho^{01}=1/2$ while $\rho^0=1$.
\end{proposition}

\begin{proof}
Write $S_m(t)=1+t+\cdots+t^{m-1}$ and use $y$ to mark ones.
Every word avoiding $0^m$ is a sequence of blocks $0^a1$,
$0\leq a<m$, followed by $0^b$, $0\leq b<m$. Thus its generating function is
\[
 F_m(t,y)=\frac{S_m(t)}{1-ytS_m(t)}.
\]
Every word avoiding $10^m$ is either entirely zero or uniquely consists
of an arbitrary initial zero string, a first $1$, and a suffix avoiding $0^m$.
Indeed, a run of at least $m$ zeros after the first $1$ has an immediately
preceding $1$ and creates the forbidden factor. The converse is immediate.
Therefore
\[
 G_m(t,y)=\frac{1+ytF_m(t,y)}{1-t}
 =\frac{1}{(1-t)(1-ytS_m(t))}.
\]

For $m\geq2$ and positive $y$ near $1$, the polynomial
$D_m(t,y)=1-y\sum_{i=1}^m t^i$ has a unique positive zero $r(y)<1$.
It is simple since $\partial_tD_m(r(y),y)<0$.
There is no other zero with modulus at most $r(y)$: the triangle inequality
would force equality in
$1\leq y\sum_{i=1}^m |t|^i\leq y\sum_{i=1}^m r(y)^i=1$,
and equality for the terms $t$ and $t^2$ forces $t=r(y)$.
Both generating functions have a simple pole there, with nonzero numerator;
the extra pole of $G_m$ at $t=1$ is farther away.

Partial fractions, as in~\cite[Theorem~IV.9, p.~256]{FS}, now give,
for each of $F_m$ and $G_m$, coefficients
$C(y)r(y)^{-n}$ plus terms exponentially smaller relative to this one,
uniformly near $y=1$, also after one $y$ derivative.
To justify the uniform assertion, choose a circle of radius strictly between
$r(1)$ and the moduli of all the other poles, subtract the simple-pole term,
and apply the coefficient bound on that circle to the remainder and its
$y$ derivative. The positive simple-pole coefficient $C(1)$ is nonzero.
It follows that both limiting frequencies are $-r'(1)/r_m$.
Differentiating $y\sum_{i=1}^m r(y)^i=1$ gives
\[
 -\frac{r'(1)}{r_m}=\frac{1}{\sum_{i=1}^m i r_m^i}.
\]
Reversal proves the equality for $0^m1$.
When $m=1$, the formulas become $F_1=(1-yt)^{-1}$ and
$G_1=((1-t)(1-yt))^{-1}$. At $y=1$ the two poles of $G_1$ coincide;
the explicit finite-word calculation in Theorem~\ref{thm:boundary}, rather
than the simple-pole argument, gives $1/2$.
\end{proof}

\begin{corollary}\label{cor:families}
Put $f_m=(\sum_{i=1}^m i r_m^i)^{-1}$ for $m\geq2$.
Besides the identities of Proposition~\ref{prop:reduction},
\[
 \rho^{01^m}=\rho^{1^m0}=\rho^{1^m}=1-f_m,
 \qquad
 \lambda^{10^m}=\lambda^{0^m1}=\lambda^{0^m}=r_m^{-1},
\]
where $\lambda^w=\lim_{n\to\infty}|\Om_n^w|^{1/n}$.
The complementary words have the same entropy. Moreover,
$f_2>f_3>\cdots>1/2$ and $f_m\to1/2$.
\end{corollary}
The limit $1/2$ is also given in~\cite[Corollary~1]{BKV};
the following argument proves strict monotonicity at the same time.
\begin{proof}
Bit-complement replaces frequency by its complement and preserves counts;
reversal preserves both counts and frequency. The dominant poles in
Proposition~\ref{prop:reduction} give the entropy identities.
The defining root equations give $1/2<r_{m+1}<r_m\leq r_2<1$.
Writing $D_m=\sum_{i=1}^m i r_m^i$, the two equations
$\sum_{i=1}^m r_m^i=\sum_{i=1}^{m+1}r_{m+1}^i=1$ yield
\[
 D_{m+1}-D_m=
 \sum_{i=1}^m(m+1-i)(r_m^i-r_{m+1}^i)>0.
\]
If $r_m\to r$, domination by $r_2^i$ gives
$\sum_{i\geq1}r^i=1$, hence $r=1/2$.
Domination by $i r_2^i$ then gives $D_m\to\sum_{i\geq1}i2^{-i}=2$.
Thus $D_m$ increases strictly to $2$, proving the assertions for $f_m$.
\end{proof}

\section{The first affirmative length}
The candidate value $\rho^{100}=f_2$ also appears in Richey's
2024 notes~\cite[Section~2.4, Eq.~(2.50)]{R}.
We compare it with the other six forbidden words explicitly.
\begin{theorem}\label{thm:three}
Among the eight forbidden binary words of length three, the unique
maximizers are $100$ and $001$, with value
$f_2=(1+1/\sqrt5)/2$. The unique minimizers are $011$ and $110$,
with value $1-f_2$. More explicitly, let $r\in(1/2,1)$ be the root of
$1-2r+r^2-r^3=0$, and put
$h=(1-r+2r^2)/(2-2r+3r^2)$. Then
\[
\begin{array}{c|cccccccc}
w&000&001&010&011&100&101&110&111\\ \hline
\rho^w&f_3&f_2&h&1-f_2&f_2&1-h&1-f_2&1-f_3
\end{array}
\]
Here $1/2<h<2/3<f_2$.
\end{theorem}
\begin{proof}
Proposition~\ref{prop:reduction} and bit-complement give all entries except
$010$ and $101$. To count words avoiding $010$, let $A,B,C$ be the generating
functions of allowable continuations when the longest suffix that is a
proper prefix of $010$ is respectively the empty word, $0$, or $01$.
Again $t$ marks length and $y$ marks ones; the empty continuation contributes
$1$ in each state. The transitions give
\[
 A=1+tyA+tB,\qquad B=1+tB+tyC,\qquad C=1+tyA.
\]
Eliminating $B,C$ gives
\[
 A(t,y)=\frac{1+t^2y}{1-t(1+y)+t^2y-t^3y^2}.
\]
At $y=1$ the denominator $P(t)=1-2t+t^2-t^3$ has derivative
$-2+2t-3t^2<0$ on the real line and one real root $r\in(1/2,1)$.
Its other two roots are a conjugate pair of modulus $r^{-1/2}>1$,
since the product of its three roots is $1$.
Thus $r$ is the simple dominant pole, the numerator is positive there,
and the pole remains separated for $y$ near $1$.
The coefficient and derivative argument from
Proposition~\ref{prop:reduction}, applied to this denominator, gives
\[
 \rho^{010}=-\frac{r'(1)}r
 =\frac{1-r+2r^2}{2-2r+3r^2}=h.
\]
The denominator $L=2-2r+3r^2$ is positive, and
\[
 h-\frac12=\frac{r^2}{2L}>0,\qquad
 \frac23-h=\frac{1-r}{3L}>0.
\]
Bit-complement gives $\rho^{101}=1-h$.
Finally $f_3<f_2$ by Corollary~\ref{cor:families}, while
$f_3,h>1/2$ and all their complements are below $1/2$.
This proves both uniqueness statements.
\end{proof}

\begin{corollary}\label{cor:alphabet}
For an alphabet of size $N\geq2$ and a fixed target symbol $b$, avoiding
the single symbol $b$ gives target frequency $0$ at every positive length.
Avoiding any other single symbol gives frequency $1/(N-1)$.
Thus the $N-1$ non-target symbols are exactly the maximizing singleton
forbidden words.
\end{corollary}
\begin{proof}
Avoiding a single symbol leaves $(N-1)^n$ words of length $n$.
If the target remains available, each of the $n$ positions is the target
in exactly $(N-1)^{n-1}$ such words. Otherwise it never occurs.
\end{proof}

\section{Coverage and verification}
The comparison not established above is whether, for every $k\geq4$,
\[
 \rho^w\leq\frac{1}{\sum_{i=1}^{k-1}i r_{k-1}^i}
 \qquad\text{for every }w\in\{0,1\}^k.
\]
Proposition~\ref{prop:reduction} computes the proposed value and
Theorem~\ref{thm:three} proves its optimality for $k=3$.
The short-length counterexamples do not decide the displayed comparison.

\begin{center}\small
\begin{tabular}{@{}p{.34\linewidth}p{.59\linewidth}@{}}
\toprule
Source item in~\cite{BMR}&Coverage\\\midrule
Section~5.3, p.~19, unrestricted extremal proposal&Refuted (explicit
counterexamples $k=1,2$). Proved for $k=3$; $k\geq4$ not covered.\\
Example~2.10, p.~5, golden-mean value&Reproved; prior proof
Lekkerkerker (1951), St.~John (1984).\\
Section~3.1, p.~8, four reducible families&Reproved for $m\geq2$;
prior proof Kova\v{c}evi\'c (2018/2019), constant-word case Chang (1986).
Corrected boundary $m=1$; complementary identities and strict decrease follow above.\\
Section~5.3, p.~19, fixed length and fixed number of ones&Length
three is classified above. The stated range $2\leq j<k/2$ is not covered.\\
Section~5.3, p.~19, fixed-border comparisons&Not covered.\\
Strict decrease of $f_m$ (not posed)&Not posed; proved in
Corollary~\ref{cor:families}. Limit previously given by Baril--Kirgizov--Vajnovszki (2021).\\
Section~6, p.~20, larger alphabets&Singleton specialization not posed; proved in
Corollary~\ref{cor:alphabet}; general symmetry questions not covered.\\
Section~6, Conjecture~6.1 (balanced borders)&Separate from the
extremal comparison; not covered.\\
\bottomrule
\end{tabular}
\end{center}

The pinned Lean declaration
\href{https://github.com/the-omega-institute/trureturing/blob/d4a6bd75fd78850fed8cf60aebb6ce7fceffa2e1/D5/S1/Words/Forbidden/BonaMagaRicheyExtremalFrequency.lean}
{\texttt{BonaMagaRicheyExtremalFrequency.result}}
states the negation of the proposal quantified over all $k\geq1$, using
the one-letter case. In that statement, Boolean \texttt{true} denotes $1$
and \texttt{false} denotes $0$. The two- and three-letter classifications,
the family identities, their strict monotonicity, and the alphabet extension
have the written proofs above.
The accompanying \texttt{verify.py} compares exact finite-word counts and
integer first moments with independently derived recurrences; its output
specifies the finite ranges. These checks support the counting identities,
whereas the infinite limits follow from the displayed proofs.

\begin{thebibliography}{8}
\bibitem{BMR}
Mikl\'os B\'ona, Bal\'azs Maga and Jacob Richey,
\emph{Letter frequency in shifts of finite type with one forbidden word},
arXiv:2606.06655v2, revised September 9, 2026.
\href{https://arxiv.org/abs/2606.06655v2}{arXiv:2606.06655v2}.
\bibitem{L}
C.~G.~Lekkerkerker,
\emph{Voorstelling van natuurlijke getallen door een som van getallen van Fibonacci},
Stichting Mathematisch Centrum, Zuivere Wiskunde ZW~30/51, 1951,
Stelling~B, p.~2.
\href{https://ir.cwi.nl/pub/6922}{CWI text}.
\bibitem{SJ}
Peter H.~St.~John,
\emph{On the asymptotic proportions of zeros and ones in Fibonacci sequences},
Fibonacci Quarterly \textbf{22}(2) (1984), 144--145.
\href{https://www.fq.math.ca/Scanned/22-2/stjohn.pdf}{Journal text}.
\bibitem{C}
Derek K.~Chang, \emph{On Fibonacci binary sequences},
Fibonacci Quarterly \textbf{24}(2) (1986), 178--179.
\href{https://www.fq.math.ca/Scanned/24-2/chang2.pdf}{Journal text}.
\bibitem{K}
Mladen Kova\v{c}evi\'c,
\emph{Runlength-limited sequences and shift-correcting codes: asymptotic analysis},
IEEE Transactions on Information Theory \textbf{65}(8) (2019), 4804--4814;
first preprint 2018, \href{https://arxiv.org/abs/1803.06117v3}{arXiv:1803.06117v3}.
\bibitem{BKV}
Jean-Luc Baril, Sergey Kirgizov and Vincent Vajnovszki,
\emph{Asymptotic bit frequency in Fibonacci words},
Pure Mathematics and Applications \textbf{30}(1) (2022), 23--30;
first preprint 2021, \href{https://arxiv.org/abs/2106.13550v2}{arXiv:2106.13550v2}.
\bibitem{R}
Jacob Richey, \emph{Pattern avoidance, entropy, and hitting time},
notes available August~12, 2024, Section~2.4, Eqs.~(2.49)--(2.54).
\href{https://github.com/jfrichey/jfrichey.github.io/blob/846c835def30a51945bac23d8ed215b61056b87c/pagedocs/sft_patterns.pdf}{Dated author text}.
\bibitem{FS}
Philippe Flajolet and Robert Sedgewick,
\emph{Analytic Combinatorics}, Cambridge University Press, 2009,
Theorem~IV.9, p.~256.
\href{https://algo.inria.fr/flajolet/Publications/book.pdf}{Author-hosted text}.
\end{thebibliography}
\medskip\noindent\small
Copyright \textcopyright\ 2026 Haobo Ma and Wenlin Zhang.
Licensed under Creative Commons Attribution 4.0 International (CC BY 4.0).
\end{document}
